\begin{description}
\item[(D2.1)]
From the table we see that the angular size of Moon is smallest close to the New Moon day.
Thus, the answer is \fbox{New Moon}. \hfill(3.0)

\emph{Justification is NOT necessary for full credit.}

\item[(D2.2)]
As there is an eclipse happening in this month, the lunar nodes are close to Full Moon day
and New Moon day. Next we notice that lowest declination of Moon is just $18^\circ$. This
means that after the New Moon day, the orbit of Moon is above the ecliptic. In other words,
the ascending node is near the \fbox{New Moon}. \hfill(4.0)

\emph{Justification is NOT necessary for full credit.}

\item[(D2.3)]
The largest angular size of the Moon in the ephemerides is $2008.3''$ and the smallest
angular size is $1763.7''$. The distance is inversely proportional to the angular size. Hence
ratio of distance at perigee to the distance at apogee is:
\begin{align*}
\mathrm{Ratio} = \frac{r_{\mathrm{perigee}}}{r_{\mathrm{apogee}}}
 &= \frac{a_0}{1+e}\times\frac{1-e}{a_0}
 & &\hspace{-2em}(2.0)\\
 % the two a_0 factors are struck through (cancelled) in the original
\therefore \frac{1-e}{1+e} &= \frac{1763.7}{2008.3} = 0.87821\\
\therefore e &= \frac{1 - 0.87821}{1 + 0.87821} = 0.064846\\
&\fbox{$e \simeq 0.065$}
 & &\hspace{-2em}(2.0)
\end{align*}
\emph{Rounding off is done to account for the fact that our data are not continuous, hence
exact angular sizes at perigee and apogee are not known. A non-rounded answer will also
receive full credit.}

\item[(D2.4)]
The following construction is shown. \hfill(5.0)
% FIGURE NEEDED: geometric construction over the eclipse composite image
% (2016_data_analysis_sol.pdf p.10): two chords AB and DC drawn across the umbra circle,
% their perpendicular bisectors intersecting at the centre O, with measured values
% D_moon = 3.3 cm and D_umbra = 9.1 cm.

Only two chords are necessary to determine the centre.

\emph{The credit is divided in two parts for the drawing:}
\begin{itemize}
\item \emph{Realisation that centre of umbra circle needs to be determined to find
$\theta_{\mathrm{umbra}}$: 1.5}
\item \emph{Accurate determination of centre of umbra circle by geometric construction: 2.5\\
Determination of centre of hand-drawn circle: maximum 1.5}
\item \emph{Measuring diameters of umbra and Moon: 0.5 each}
\end{itemize}

By estimating approximate centre of the shadow in the image, we find out,
\begin{align*}
\frac{\theta_{\mathrm{umbra}}}{\theta_{\mathrm{Moon}}}
 &= \frac{d_{\mathrm{umbra}}}{d_{\mathrm{Moon}}} & &\hspace{-2em}(2.0)\\
 &= \frac{9.1}{3.3} = 2.76\\
\therefore \theta_{\mathrm{umbra}} &\simeq 2.76\,\theta_{\mathrm{Moon}}
 & &\hspace{-2em}(1.0)
\end{align*}
\emph{Acceptable range: $\pm0.10$.}

\item[(D2.5)]
The following diagram needs to be drawn.
% FIGURE NEEDED: ray diagram (2016_data_analysis_sol.pdf p.11): Sun's rays S1 and S2
% converging past the Earth (centre O, radius R_earth, terminator points E, P), crossing at
% J on the axis; the Moon's path through points Q, B, A; angles theta_1, theta_2 and
% theta_Sun marked.

Angular size of umbra is $\theta_{\mathrm{umbra}} = 2\angle BOQ$\\
Angular size of penumbra is $\theta_{\mathrm{penumbra}} = 2\angle AOQ$

We have
\begin{align*}
QA &= QP + PA\\
 &= OE + PE\tan\theta_1\\
 &\approx R_\oplus + d_{\mathrm{Moon}}\theta_1 & &\text{since } PA \ll PE\\
 &\approx R_\oplus + d_{\mathrm{Moon}}\frac{\theta_{\mathrm{Sun}}}{2}
   & &\text{since } \theta_1 \approx \theta_2 \approx \theta_{\mathrm{Sun}}/2
   \quad(2.0)
\end{align*}
and
\begin{align*}
QB &= QP - PB\\
 &= OE - PE\tan\theta_2\\
 &\approx R_\oplus - d_{\mathrm{Moon}}\theta_2\\
 &\approx R_\oplus - d_{\mathrm{Moon}}\frac{\theta_{\mathrm{Sun}}}{2}
   & &\hspace{-2em}(2.0)
\end{align*}
\begin{align*}
\therefore \theta_{\mathrm{penumbra}} = 2\angle AOQ
 = 2\tan^{-1}\left(\frac{QA}{OQ}\right) &\approx 2\frac{QA}{OQ}
   \qquad\text{since } QA \ll OQ\\
 = 2\frac{R_\oplus + d_{\mathrm{Moon}}\dfrac{\theta_{\mathrm{Sun}}}{2}}{d_{\mathrm{Moon}}}
 &= \frac{2R_\oplus}{d_{\mathrm{Moon}}} + \theta_{\mathrm{Sun}}
   & &\hspace{-2em}(1.0)
\end{align*}
and
\begin{align*}
\theta_{\mathrm{umbra}} = 2\angle BOQ
 = 2\tan^{-1}\left(\frac{QB}{OQ}\right) &\approx 2\frac{QB}{OQ}\\
 = 2\frac{R_\oplus - d_{\mathrm{Moon}}\dfrac{\theta_{\mathrm{Sun}}}{2}}{d_{\mathrm{Moon}}}
 &= \frac{2R_\oplus}{d_{\mathrm{Moon}}} - \theta_{\mathrm{Sun}}
   & &\hspace{-2em}(1.0)
\end{align*}
Subtracting,
\[
\theta_{\mathrm{penumbra}} - \theta_{\mathrm{umbra}} = 2\theta_{\mathrm{Sun}}
\Rightarrow \theta_{\mathrm{penumbra}} = \theta_{\mathrm{umbra}} + 2\theta_{\mathrm{Sun}}
\]
\hfill(1.0)

We have,
\[
\theta_{\mathrm{umbra}} = 2.76\,\theta_{\mathrm{Moon}}
\qquad\text{and } \theta_{\mathrm{Sun}} = 1915.0''
\]
From the given data, $\theta_{\mathrm{Moon}} = 2008.3''$. \hfill(1.0)

Therefore,
\begin{align*}
\theta_{\mathrm{penumbra}} &= 2.76\,\theta_{\mathrm{Moon}}
 + 2\frac{1915.0}{2008.3}\theta_{\mathrm{Moon}}\\
&\fbox{$\theta_{\mathrm{penumbra}} = 4.67\,\theta_{\mathrm{Moon}}$}
 & &\hspace{-2em}(1.0)
\end{align*}
\emph{Acceptable range: $4.57\,\theta_{\mathrm{Moon}}$ to $4.77\,\theta_{\mathrm{Moon}}$.}

\underline{Alternative solution}:

In the figure below, rays $HEA$ and $IFC$ are coming from one edge of solar disk and rays
$HFD$ and $GEB$ are coming from the opposite edge. The observer ($O$) is assumed to be at
the centre of the Earth. The Moon travels along the path $ABCD$ during the course of
eclipse.
% FIGURE NEEDED: two ray diagrams (2016_data_analysis_sol.pdf p.13): Earth with tangent
% rays from both solar limbs crossing at J, the Moon's path ABCD, points E, F on the
% Earth's limb, G, H, I beyond it, and the equal angles theta_Sun marked at several
% vertices.

From figure,
\[
\angle AEB = \angle GEH = \angle HFI = \angle DFC = \angle EJF = \theta_{\mathrm{Sun}}
\]
\hfill(3.0)
\begin{align*}
\theta_{\mathrm{umbra}} &= \angle BOC = 2.76\,\theta_{\mathrm{Moon}}
 & &\hspace{-2em}(1.0)\\
\theta_{\mathrm{penumbra}} &= \angle AOD & &\hspace{-2em}(1.0)\\
\angle AOD &= \angle AOB + \angle BOC + \angle COD\\
\angle AOB &= \angle AEB & &\hspace{-2em}(1.0)\\
\angle COD &= \angle CFD & &\hspace{-2em}(1.0)\\
\theta_{\mathrm{penumbra}} &\simeq \angle AEB + \theta_{\mathrm{umbra}} + \angle CFD\\
 &= 2\theta_{\mathrm{Sun}} + 2.76\,\theta_{\mathrm{Moon}}
 = 2\times1915.0'' + 2.76\times2008.3''\\
\theta_{\mathrm{penumbra}} &= 9372.9'' = 4.67\,\theta_{\mathrm{Moon}}
 & &\hspace{-2em}(2.0)
\end{align*}

\item[(D2.6)]
From the Moon,
\[
\theta_{\mathrm{Earth}} = \frac{2R_\oplus}{d_{\mathrm{Moon}}}
\]
\hfill(1.0)

From part (D2.5),
\[
\theta_{\mathrm{umbra}} + \theta_{\mathrm{penumbra}} = 2\theta_{\mathrm{Earth}}
\]
\hfill(2.0)
\begin{align*}
\therefore \theta_{\mathrm{Earth}}
 &= \frac{\theta_{\mathrm{umbra}} + \theta_{\mathrm{penumbra}}}{2}
 = \frac{2.76 + 4.67}{2}\theta_{\mathrm{Moon}} = 3.72\,\theta_{\mathrm{Moon}}\\
&\fbox{$\theta_{\mathrm{Moon}} = {} = 0.269\,\theta_{\mathrm{Earth}}$}
 & &\hspace{-2em}(2.0)
\end{align*}
% sic: the boxed answer prints a doubled equals sign.

\underline{Alternative solution}:

Let us say that the Moon is at position of $B$. Thus, angular size of Earth as seen from
this position will be, (see figure in the previous part)
\begin{align*}
\theta_{\mathrm{Earth}} &= \angle EBF = \angle BFD & &\hspace{-2em}(2.0)\\
 &= \angle BFC + \angle CFD\\
 &\simeq \theta_{\mathrm{umbra}} + \theta_{\mathrm{Sun}} & &\hspace{-2em}(1.0)
\end{align*}
The angular size of the Full Moon on 28 September as seen in the table is $2008.3''$.
\begin{align*}
\theta_{\mathrm{Earth}} &= 2.76\times2008.3'' + 1915.0'' = 7453.0''\\
&\fbox{$\theta_{\mathrm{Moon}} = 0.269\,\theta_{\mathrm{Earth}}
 = \dfrac{\theta_{\mathrm{Earth}}}{3.72}$}
 & &\hspace{-2em}(2.0)
\end{align*}

\item[(D2.7)]
Thus, the radius of Moon will be,
\begin{align*}
R_{\mathrm{Moon}} &= \frac{R_\oplus}{3.72} & &\hspace{-2em}(1.0)\\
R_{\mathrm{Moon}} &= \frac{6371}{3.72}\\
&\fbox{$R_{\mathrm{Moon}} \simeq 1713\,\mathrm{km}$} & &\hspace{-2em}(2.0)
\end{align*}
\emph{Acceptable range: $\pm20\,\mathrm{km}$.}

\item[(D2.8)]
The shortest and longest distances will be,
\begin{align*}
r_{\mathrm{perigee}} &= \frac{2\times1713\times206265}{2008.3}\\
&\fbox{$r_{\mathrm{perigee}} = 3.52\times10^5\,\mathrm{km}$} & &\hspace{-2em}(2.0)\\
r_{\mathrm{apogee}} &= \frac{2\times1713\times206265}{1763.7}\\
&\fbox{$r_{\mathrm{apogee}} = 4.01\times10^5\,\mathrm{km}$} & &\hspace{-2em}(2.0)
\end{align*}

\item[(D2.9)]
% FIGURE NEEDED: triangle diagram (2016_data_analysis_sol.pdf p.15): Earth (E), Moon (M) on
% its orbit and Sun (S), with the angles phi_E at Earth, phi_M at Moon and phi_S at Sun.
On September 10, phase of Moon is 0.097 and elongation of Moon is $36.2^\circ$. Angular
size of the Moon on this day is $1792.0''$. Therefore, distance to Moon (from Earth) on
September 10 is
\begin{align*}
d_{\mathrm{Moon,10}} &= \frac{2\times1713\times206265}{1792.0}\\
 &= 3.94\times10^5\,\mathrm{km} & &\hspace{-2em}(2.0)
\end{align*}
Let
\begin{align*}
\angle EMS &= \varphi_M\\
\angle ESM &= \varphi_S\\
\angle SEM &= \varphi_E\\
\therefore \varphi_E &= 36.2^\circ & &\hspace{-2em}(1.0)\\
\mathrm{phase} &= \frac{1 + \cos\varphi_M}{2} & &\hspace{-2em}(2.0)\\
\therefore \varphi_M &= \cos^{-1}\left(2\times\mathrm{phase} - 1\right)\\
 &= \cos^{-1}\left(2\times0.097 - 1\right) = \cos^{-1}(-0.806)\\
 &= 143.71^\circ & &\hspace{-2em}(2.0)\\
\varphi_S &= 180^\circ - \varphi_E - \varphi_M\\
 &= 180^\circ - 36.2^\circ - 143.71^\circ\\
 &= 0.09^\circ & &\hspace{-2em}(1.0)
\end{align*}
Now using sine rule,
\begin{align*}
\frac{d_{\mathrm{Sun}}}{d_{\mathrm{Moon,10}}} &= \frac{\sin\varphi_M}{\sin\varphi_S}
 & &\hspace{-2em}(2.0)\\
\therefore d_{\mathrm{Sun}} &= 3.94\times10^8\times\frac{\sin 143.71^\circ}{\sin 0.09^\circ}\\
&\fbox{$d_{\mathrm{Sun}} = 1.48\times10^{11}\,\mathrm{m}$} & &\hspace{-2em}(1.0)
\end{align*}
\end{description}
